\scanpage{11}
% source lines 1-3, continuation of 3 degree
\uline{exactement} \inserted{(contient trivialement)} \unsure{l'im.} \deleted{des classes} $N_{L/K}C(L)$. En d'autres termes le noyau de ${\unsure{D|\mathscr D}}_{\underline m}(K)\unsure{-|\longrightarrow}G/G'$ est \uline{exactement} \inserted{(contient trivialement)} $P_{\underline m}(K).N_{L/K}(D_{\underline m'}(L))\deleted{\unsure{\}}}$.

% source lines 4-6
Cela montre déjà que les extensions \uline{abéliennes} \inserted{\uline{finies}} de $K^*$ correspondent à \deleted{\unsure{des}} \uline{certains} sous-groupes ouverts de $C(K)$. Il y a de plus un

% source lines 7-15, circled 4 degree
\textcircled{$4^\circ$} \uline{Théorème d'existence} : \deleted{Pour} Tout sous-groupe ouvert de $C(K)$ correspond à une extension abélienne, i.e. Pour \unsure{tout|$\forall$} module $\underline m$, il existe une extension abélienne $L$ de $K$ \inserted{admettant le module $\underline m$ et} telle que
\[
N_{L/K}(D_{\underline m}(L))\subset\deleted{\unsure{D_{\underline m}}}P_{\underline m}(K).
\]
[ou aussi : pour $\forall$ ss-groupe ouvert $U$ de $C(K)$, existe une extension abélienne $L$ de $K$, t.q. $N_{L/K}C(L)=U$].

% source lines 16-21
N.B. \unsure{L'essentiel} de ces assertions ne concerne que les extensions abéliennes \deleted{\unsure{les autres}} (en fait, la sous-extension abélienne maximale $L_{\unsure{a|0|o}}$ de $K$ dans $L$). Les assertions faites équivalent \unsure{à|aux} ces assertions pour $L$ abélien, plus le fait que $N_{L/K}C(L)=N_{L_a/K}C(L_a)$
% archive sheet number: 10
